\documentclass[11pt]{article}
\usepackage[a4paper,margin=1in]{geometry}
\usepackage{amsmath,amsthm,amssymb}
\usepackage{mathtools}

\newtheorem{theorem}{Theorem}[section]
\newtheorem{proposition}[theorem]{Proposition}
\newtheorem{lemma}[theorem]{Lemma}
\newtheorem{corollary}[theorem]{Corollary}
\theoremstyle{definition}
\newtheorem{definition}[theorem]{Definition}
\newtheorem{remark}[theorem]{Remark}

\newcommand{\R}{\mathbb{R}}
\newcommand{\Fr}{\mathcal{A}}
\newcommand{\Om}{\Omega}
\newcommand{\eps}{\varepsilon}
\newcommand{\rstar}{\rho_{*}}
\newcommand{\rhad}{\widehat\rho}
\newcommand{\X}{\mathcal X}

\title{Boundary-Layer Transfer and Bounded Sharp Saint--Venant\\
       (Route~$\delta$, Plan~2, Building Block \textsc{BoundaryLayerTransfer})}
\author{}
\date{}

\begin{document}
\maketitle

\begin{abstract}
We close the route-$\delta$ chain \emph{relative to the upstream metric
trace input} by transferring the metric-trace estimate at the deficit
endpoint $\rhad:=\rho_\delta=1-\kappa\sqrt{\delta_T}$ back to the
original set $\Om$.
The upstream weighted metric trace (\cite{WeightedMetricTrace}) supplies
the uniform pointwise bound
$\sup_{\rho\in J^*}d_{\X}(F_\rho,B_1)^2\le C_*\,\delta_T(\Om)$ on the
order-one window $J^*=[\rho_\delta-L^*,\rho_\delta]$; specialising to
the upper endpoint gives $d_{\X}(F_{\rho_\delta},B_1)^2\le C_*\,\delta_T(\Om)$.
The boundary-layer volume
$|\Om\setminus E_{\rho_\delta}|=\omega_n(1-\rho_\delta^{\,n})$ is exactly
$O(\sqrt{\delta_T})$ by the definition of $\rho_\delta$; the elementary
superlevel-transfer estimate of \cite{LevelSetIdentity} (Lemma~8.3)
gives $\Fr(\Om)\le \Fr(E_{\rho_\delta})+C_{\rm lay}\sqrt{\delta_T}$.
Squaring
with $(a+b)^2\le 2a^2+2b^2$ doubles the constant on $\delta_T$ but
preserves the sharp exponent~$2$. We obtain the bounded sharp
Saint--Venant inequality
$\Fr(\Om)^2\le C(n,R,\rstar)\,\delta_T(\Om)$ for all open
$\Om\subset B_R$ with $|\Om|=\omega_n$ inside an explicit smallness
window $\delta_T(\Om)\le \delta_0(n,R,\rstar)$. Outside that window
the inequality is trivial after enlarging $C$ by the factor
$4/\delta_0$. The small-deficit constant is polynomial in
$1/\rstar$, $R$, $n$ and the constants of the upstream building
blocks. The full bounded constant has this same polynomial dependence
provided the inverse upstream smallness threshold $\delta_0^{-1}$ is
polynomially controlled.
\end{abstract}

\section{Statement of the theorem}\label{sec:statement}

Throughout, $n\ge 2$, $R\ge 2$, $\rstar\in(0,1)$ are fixed.
$\Om\subset\R^n$ denotes a bounded open set with $|\Om|=\omega_n$ and
$\Om\subset B_R$. The torsion functional is normalised by
$\delta_T(\Om):=E(\Om)-E(B_1)\ge 0$.
The Fraenkel asymmetry (cf.\ \cite{AFP2000}) is
\[
\Fr(\Om):=\inf\!\left\{\frac{|\Om\Delta(z+B_1)|}{\omega_n}\;:\;z\in\R^n\right\}\in[0,2).
\]
For $\rho\in[\rstar,1]$ we write $E_\rho=\{u_\Om>t(\rho)\}$ where
$|E_\rho|=\omega_n\rho^n$ and $u_\Om$ is the torsion function of $\Om$
(see \cite{LevelSetIdentity}). The normalised level set
$F_\rho:=\rho^{-1}E_\rho\in\X$, where $\X$ is the metric quotient of
finite-measure sets by translations equipped with
\[
d_{\X}([A],[C]):=\inf_{a\in\R^n}|A\Delta(C+a)|.
\]

\begin{theorem}[Bounded sharp Saint--Venant; building block
\textsc{BoundaryLayerTransfer}]\label{thm:main}
There exist
$C=C(n,R,\rstar)>0$ and $\delta_0=\delta_0(n,R,\rstar)>0$ such that
every open $\Om\subset B_R$ with $|\Om|=\omega_n$ satisfies
\begin{equation}\label{eq:main}
\Fr(\Om)^{2}\le C(n,R,\rstar)\,\delta_T(\Om).
\end{equation}
The small-deficit constant is polynomial in $1/\rstar$, $R$, $n$
and the constants of \cite{WeightedMetricTrace}; the final all-deficit
constant also contains $4/\delta_0(n,R,\rstar)$. Thus the advertised
polynomial dependence for $C(n,R,\rstar)$ is conditional on polynomial
control of $\delta_0^{-1}$ from the upstream metric-trace package.
The exponent $2$ is sharp.
We further assume $\Om$ is open of full measure $|\Om|=\omega_n$
contained in $B_R$, sufficient that the torsion function
$u_\Om\in W^{1,2}_0(\Om)$ exists with level sets $E_\rho$ of finite
perimeter for a.e.\ $\rho$ (\cite{Maggi2012} coarea).
\end{theorem}

The proof occupies \S\ref{sec:setup}--\S\ref{sec:smallness}: the
small-deficit regime $\delta_T(\Om)\le \delta_0$ is closed by the
metric trace, the elementary superlevel transfer, and the
squaring identity; the complementary regime $\delta_T(\Om)\ge\delta_0$
is closed by the trivial bound $\Fr(\Om)\le 2$ after absorbing
$4/\delta_0$ into $C$.

\section{The metric-trace input}\label{sec:setup}

We record, in the form needed below, the output of the weighted
metric trace lemma assembled in \cite{WeightedMetricTrace}. Set
\begin{equation}\label{eq:rho-delta}
\rho_\delta:=1-\kappa\sqrt{\delta_T(\Om)},
\end{equation}
where $\kappa=\kappa(n,\rstar)>0$ is the profile-gap constant fixed
in \cite{LevelSetIdentity}; equivalently
$\rho_\delta$ is the largest radius for which the profile-gap lower
bound on the measure $d\mu(\rho)=(-t'(\rho))\,d\rho$ holds (cf.\
\cite[(5.1)]{LevelSetIdentity}).

\begin{remark}[No pointwise lower-gradient input]\label{rem:hypG}
The constant $C_*=C_*(n,R,\rstar)$ inherited from
\cite[Thm.~1.3]{WeightedMetricTrace} no longer depends on a pointwise
lower bound for $|\nabla u|$ on level surfaces.  The former Hyp-G input
has been replaced upstream by the self-truncated velocity estimate of
\cite{TrimmedVelocityRepair}: on Fusco--Julin good levels, the
low-gradient part of $\partial^*E_\rho$ is charged by the
Cauchy--Schwarz defect $D_H$ itself.  The boundary-layer transfer
arguments of the present block do not use any gradient lower bound
directly.
\end{remark}

Downstream, \textsc{GlobalAssembly} will specialize this bounded
statement to the dimensional radius $R_*(n)$ produced by the bounded
reduction \cite{BDV-BoundedReduction-Plan1}.

\begin{theorem}[Metric trace at $\rho_\delta$; upstream input from \cite{WeightedMetricTrace}]
\label{thm:metric-trace}
There exist $C_*=C_*(n,R,\rstar)>0$ and
$\delta_1=\delta_1(n,R,\rstar)>0$ such that for every open
$\Om\subset B_R$ with $|\Om|=\omega_n$ and
$\delta_T(\Om)\le\delta_1$, setting
\begin{equation}\label{eq:rhad-window}
\rhad\;:=\;\rho_\delta\;=\;1-\kappa\sqrt{\delta_T(\Om)},
\end{equation}
we have
\begin{equation}\label{eq:metric-trace}
d_{\X}(F_{\rhad},B_1)^{2}\le C_*\,\delta_T(\Om).
\end{equation}
In fact the bound holds uniformly for $\rho\in J^*=[\rho_\delta-L^*,\rho_\delta]$
with $L^*=(\rho_\delta-\rstar)/4$.
\end{theorem}

In the sequel we always impose the additional admissibility threshold
\begin{equation}\label{eq:delta-adm}
  \delta_{\rm adm}(n,\rstar)
  :=\left(\frac{1-\rstar}{2\kappa(n,\rstar)}\right)^2 .
\end{equation}
If $\delta_T(\Om)\le\delta_{\rm adm}$, then
\begin{equation}\label{eq:rho-adm}
  \rho_\delta\ge\frac{1+\rstar}{2}>\rstar,
  \qquad
  L^*=\frac{\rho_\delta-\rstar}{4}\ge\frac{1-\rstar}{8}.
\end{equation}
Thus the endpoint $\rho_\delta$ is admissible and the trace window
$J^*$ is nonempty with length bounded below by an order-one collar
constant depending only on $(n,\rstar)$.

\paragraph{Reference.}
This is the conclusion of \cite[Thm.~5.6, Cor.~5.7]{WeightedMetricTrace},
combining the integrated metric-derivative bound
\cite[(4.2)]{WeightedMetricTrace} (which encodes the integrated
Cauchy--Schwarz defect $D_H$) with the profile-gap measure
\cite[(5.1)]{LevelSetIdentity} and the Fusco--Julin
isoperimetric input \cite[Lemma~3.1]{WeightedMetricTrace}
(which encodes the integrated isoperimetric defect $D_I$).  The bound
at the endpoint $\rho_\delta$ follows from Markov on the distance term
plus kinetic Cauchy--Schwarz transport across $J^*$; see
\cite[Rem.~5.7]{WeightedMetricTrace} for the structural exposition.
No additional input is needed here.

\begin{remark}\label{rem:asym-as-distance}
By the definition of $d_{\X}$ and the volume normalisation
$|E_{\rhad}|=\omega_n\rhad^{\,n}$, there exists
$z=z(\rhad)\in\R^n$ such that
$|E_{\rhad}\Delta(z+B_{\rhad})|\le \rhad^{\,n}d_{\X}(F_{\rhad},B_1)$.
Dividing by $|E_{\rhad}|=\omega_n\rhad^{\,n}$,
\begin{equation}\label{eq:Aerho}
\Fr(E_{\rhad})\le \omega_n^{-1}d_{\X}(F_{\rhad},B_1),
\qquad
\Fr(E_{\rhad})^{2}\le \omega_n^{-2}C_*\delta_T(\Om).
\end{equation}
Here $\Fr(E_{\rhad})$ is the Fraenkel asymmetry of $E_{\rhad}$
relative to balls of \emph{its own volume} $\omega_n\rhad^{\,n}$,
not of $\omega_n$.
\end{remark}

\section{Volume of the removed boundary layer}\label{sec:layer-volume}

\begin{lemma}[Boundary-layer volume]\label{lem:layer-vol}
Let $\rhad=\rho_\delta$ as in Theorem~\ref{thm:metric-trace}. There is
$C_{\rm lay}=C_{\rm lay}(n,\rstar)>0$ such that, for
$\delta_T(\Om)\le\delta_2(n,R,\rstar)$ small,
\begin{equation}\label{eq:layer-vol}
|\Om\setminus E_{\rhad}|\;=\;\omega_n\bigl(1-\rhad^{\,n}\bigr)
\;\le\; C_{\rm lay}\sqrt{\delta_T(\Om)}.
\end{equation}
\end{lemma}

\begin{proof}
By \eqref{eq:rho-delta} and \eqref{eq:rhad-window},
$\rhad=1-\kappa\sqrt{\delta_T(\Om)}$.  The map $\rho\mapsto 1-\rho^n$ is
$C^1$ with derivative $-n\rho^{n-1}$, so the mean-value theorem applied
on $[\rhad,1]\subset[\rstar,1]$ gives
\[
   1-\rhad^{\,n}\;\le\;n(1-\rhad)\;=\;n\kappa\sqrt{\delta_T(\Om)}.
\]
This proves \eqref{eq:layer-vol} with
$C_{\rm lay}:=\omega_n n\kappa$. Equivalently, $C'_n:=n\kappa$ so that
$C_{\rm lay}=\omega_n C'_n$ and the
prefactor in \eqref{eq:transfer-rhat} simplifies to $2C'_n$ without
double-cancellation of $\omega_n$.  Choosing
$\delta_2:=\min\{\delta_1,\delta_{\rm adm}\}$ ensures both
\eqref{eq:rhad-window} and the admissibility bounds
\eqref{eq:rho-adm} apply.
\end{proof}

\begin{remark}[Why $\rhad=\rho_\delta$ is the natural choice]
\label{rem:why-rho-delta}
The upstream weighted metric trace
\cite[Thm.~5.6, Rem.~5.7]{WeightedMetricTrace} delivers a uniform
pointwise bound $d_\X(F_\rho,B_1)^2\le C_*\delta_T$ valid for
\emph{every} $\rho$ in the order-one window $J^*=[\rho_\delta-L^*,\rho_\delta]$.
The boundary-layer transfer cost is monotone in $\rhad$: smaller
$\rhad$ leaves a larger boundary layer
$|\Om\setminus E_{\rhad}|\asymp 1-\rhad$.  Since both costs --- the
metric-trace bound and the layer-volume bound --- balance optimally at
$\rhad=\rho_\delta$, where $1-\rhad=\kappa\sqrt{\delta_T}$ is the
smallest layer in $J^*$, the endpoint choice is the unique sharp one.
\end{remark}

\section{The asymmetry transfer lemma}\label{sec:transfer}

We now prove the elementary superlevel transfer inequality
(\cite{LevelSetIdentity}, Lemma~8.3). The point is that the
asymmetry of $\Om$ at scale $|\Om|=\omega_n$ is dominated by the
asymmetry of any nested superlevel $E_t$ at its own scale
$|E_t|=\omega_n\rho^n$ plus a volume-defect correction; the
ball-scale correction $|B_{\omega_n}\setminus B_s|$ is controlled by
the same defect by elementary monotonicity.

Throughout this section, $\Om$ has $|\Om|=L=\omega_n$, $E_t\subset\Om$
is open with $|E_t|=s$ and $\eta:=L-s\in[0,L]$. We write $B_r$ for the
open ball of radius $r$ centred at the origin and recall that the
volume of $B_r$ is $\omega_n r^n$, so the ball of volume $s$ is
$B_{r_s}$ with $r_s:=(s/\omega_n)^{1/n}$, and the ball of volume $L$ is
$B_1$.

\begin{lemma}[Superlevel asymmetry transfer; cf.\
\cite{LevelSetIdentity} Lemma~8.3]\label{lem:8.3}
Let $\Om,E_t$ be as above with $E_t\subset\Om$ and $\eta=|\Om|-|E_t|$.
Then
\begin{equation}\label{eq:8.3}
\Fr(\Om)\le \Fr(E_t)+\frac{2\eta}{L}.
\end{equation}
Equivalently, since $L=\omega_n$,
\begin{equation}\label{eq:8.3-omega}
\Fr(\Om)\le \Fr(E_t)+\frac{2}{\omega_n}\,|\Om\setminus E_t|.
\end{equation}
\end{lemma}

\begin{proof}
Pick $\eps\in(0,L)$ and choose $z\in\R^n$ such that
\begin{equation}\label{eq:near-opt}
|E_t\Delta(z+B_{r_s})|\le s\,\bigl(\Fr(E_t)+\eps\bigr).
\end{equation}
We compare $\Om$ to the ball $z+B_1$ of volume $L$. The triangle
inequality for symmetric differences gives
\begin{equation}\label{eq:tri}
|\Om\Delta(z+B_1)|\le |\Om\Delta E_t|+|E_t\Delta(z+B_{r_s})|+|(z+B_{r_s})\Delta(z+B_1)|.
\end{equation}
We bound each summand.

\emph{First summand.} Because $E_t\subset\Om$,
$\Om\Delta E_t=\Om\setminus E_t$, hence
\begin{equation}\label{eq:term1}
|\Om\Delta E_t|=\eta.
\end{equation}

\emph{Second summand.} \eqref{eq:near-opt} gives
\begin{equation}\label{eq:term2}
|E_t\Delta(z+B_{r_s})|\le s\bigl(\Fr(E_t)+\eps\bigr).
\end{equation}

\emph{Third summand.} Since $z+B_{r_s}\subset z+B_1$,
\begin{equation}\label{eq:term3}
|(z+B_{r_s})\Delta(z+B_1)|=|B_1\setminus B_{r_s}|=L-s=\eta.
\end{equation}

Combining \eqref{eq:tri}--\eqref{eq:term3} and dividing by $L$,
\begin{equation}\label{eq:combine}
\frac{|\Om\Delta(z+B_1)|}{L}
\le \frac{2\eta}{L}+\frac{s}{L}\bigl(\Fr(E_t)+\eps\bigr)
\le \frac{2\eta}{L}+\Fr(E_t)+\eps,
\end{equation}
using $s\le L$. Since $\Fr(\Om)\le L^{-1}|\Om\Delta(z+B_1)|$ and
$\eps\in(0,L)$ is arbitrary, \eqref{eq:8.3} follows.
\end{proof}

\begin{remark}[On the ball-scale correction]
\label{rem:ball-scale}
The triangle inequality \eqref{eq:tri} contains two distinct
volume-defect terms: $|\Om\Delta E_t|=\eta$ (the boundary layer of
$\Om$ outside $E_t$) and $|B_1\setminus B_{r_s}|=\eta$ (the
``ball-scale'' correction). They add to give $2\eta/L$. There is no
$|B_s|/|B_L|$ rescaling to worry about, because the symmetric
difference $|E_t\Delta(z+B_{r_s})|$ is measured at the natural scale
$s$ of $E_t$, and the cross term to scale $L$ is precisely
$|B_1\setminus B_{r_s}|$. This is the same combinatorial identity
used in \cite[Lemma~8.3]{LevelSetIdentity} and in the
$L^1$-rearrangement bookkeeping in \cite{Maggi2012}
(see also \cite{FigalliMaggi2011}).
\end{remark}

\begin{corollary}[Transfer at the level $\rhad$]\label{cor:transfer}
Let $\Om\subset B_R$ with $|\Om|=\omega_n$,
$\delta_T(\Om)\le\delta_2$, and let $\rhad$ be as in
Theorem~\ref{thm:metric-trace}. Then
\begin{equation}\label{eq:transfer-rhat}
\Fr(\Om)\le \Fr(E_{\rhad})+2C'_n\sqrt{\delta_T(\Om)},
\end{equation}
where $C'_n=n\kappa$ (equivalently
$C_{\rm lay}=\omega_n C'_n$) is the boundary-layer constant of
Lemma~\ref{lem:layer-vol}.
\end{corollary}

\begin{proof}
$E_{\rho_\delta}\subset\Om$ by definition, and
$|\Om\setminus E_{\rho_\delta}|\le C_{\rm lay}\sqrt{\delta_T(\Om)}$ by
Lemma~\ref{lem:layer-vol}. Apply \eqref{eq:8.3-omega} with
$E_t=E_{\rho_\delta}$.
\end{proof}

\section{Squaring preserves the exponent~$2$}\label{sec:squaring}

We combine Corollary~\ref{cor:transfer} with
\eqref{eq:Aerho} and the elementary inequality
$(a+b)^2\le 2a^2+2b^2$.

\begin{theorem}[Small-$\delta_T$ bound]\label{thm:small-delta}
There exist $C_3=C_3(n,R,\rstar)>0$ and
$\delta_0=\delta_0(n,R,\rstar)\in(0,\delta_2]$ such that every open
$\Om\subset B_R$ with $|\Om|=\omega_n$ and
$\delta_T(\Om)\le \delta_0$ satisfies
\begin{equation}\label{eq:small-delta}
\Fr(\Om)^{2}\le C_3(n,R,\rstar)\,\delta_T(\Om).
\end{equation}
\end{theorem}

\begin{proof}
Take $\delta_T(\Om)\le\delta_2$ so that
Corollary~\ref{cor:transfer} applies. Set
$a:=\Fr(E_{\rhad})$, $b:=2C'_n\sqrt{\delta_T(\Om)}$.
Then \eqref{eq:transfer-rhat} reads $\Fr(\Om)\le a+b$. Squaring,
\begin{equation}\label{eq:square}
\Fr(\Om)^2\le 2a^2+2b^2
=2\Fr(E_{\rhad})^2+8(C'_n)^{2}\,\delta_T(\Om).
\end{equation}
By \eqref{eq:Aerho},
$\Fr(E_{\rhad})^2\le \omega_n^{-2}C_*\delta_T(\Om)$. Substituting,
\begin{equation}\label{eq:final-const}
\Fr(\Om)^2\le \left(\frac{2C_*}{\omega_n^{2}}+8(C'_n)^{2}\right)\delta_T(\Om)
=:C_3\,\delta_T(\Om).
\end{equation}
This proves \eqref{eq:small-delta} with
$\delta_0:=\delta_2=\min\{\delta_1,\delta_{\rm adm}\}$. The
exponent~$2$ is preserved because $(a+b)^2\le 2a^2+2b^2$
\emph{doubles the constant} but does not raise either summand to a
higher power.
\end{proof}

\begin{remark}[The factor of~$2$]\label{rem:factor-2}
The identity $(a+b)^2=a^2+2ab+b^2\le 2(a^2+b^2)$, valid by
$2ab\le a^2+b^2$, is sharp up to the factor of~$2$ (saturated when
$a=b$). Replacing it by Young's inequality with a parameter
$\theta\in(0,1)$, $(a+b)^2\le (1+\theta^{-1})a^2+(1+\theta)b^2$,
allows the prefactor on $\Fr(E_{\rhad})^2$ to be made arbitrarily
close to~$1$, at the cost of inflating the coefficient on
$\delta_T$. This does not affect the exponent, only the constant
$C_3$. For the bounded sharp Saint--Venant inequality the cleanest
choice is $\theta=1$, giving \eqref{eq:final-const}.
\end{remark}

\section{Constants chain}\label{sec:constants}

The constant $C_3=C_3(n,R,\rstar)$ produced by
Theorem~\ref{thm:small-delta} unwinds as
\begin{equation}\label{eq:C3-chain}
C_3=\frac{2C_*}{\omega_n^{2}}+8(C'_n)^{2},
\qquad C'_n=n\kappa,
\end{equation}
where the inputs are:

\begin{itemize}
\item $C_*=C_*(n,R,\rstar)$: the metric trace constant of
Theorem~\ref{thm:metric-trace}, supplied by
\cite[Theorem on weighted metric trace]{WeightedMetricTrace}. In the
repaired metric-trace block, $C_*$ is polynomial in
$(1/\rstar, R, n)$ and in the constants of the Fusco--Julin and
oriented radial-trace inputs.

\item $C_{\rm lay}=\omega_n n\kappa$ from Lemma~\ref{lem:layer-vol},
where $\kappa=\kappa(n,\rstar)$ is the profile-gap constant of
\cite{LevelSetIdentity}.  The simpler form (no $C_*$ contribution)
reflects the endpoint choice $\rhad=\rho_\delta$, see
Remark~\ref{rem:why-rho-delta}.

\item $\omega_n=|B_1|$ is purely dimensional.
\end{itemize}

In particular, $C_3$ is polynomial in $(1/\rstar, R, n)$ and in the
explicit constants of the repaired metric trace. The full constant in
Theorem~\ref{thm:main} additionally contains
$4/\delta_0$, with
\[
  \delta_0=\min\left\{\delta_1(n,R,\rstar),
  \left(\frac{1-\rstar}{2\kappa(n,\rstar)}\right)^2\right\}.
\]
Consequently the final polynomial-dependence claim is valid once the
upstream metric trace supplies polynomial control of
$\delta_1^{-1}$. Specialising to $R=R_*(n)$ (the bounded reduction
radius) and $\rstar=1/2$ as in \cite{WeightedMetricTrace}, the
constant is dimensional.

\begin{remark}[Sharpness of the dependence]
The polynomial dependence on $1/\rstar$ enters through (i) the
profile-gap constant in \cite[(5.1)]{LevelSetIdentity}; (ii) the
Cauchy--Schwarz bookkeeping in the metric derivative bound
\cite[(2.4),(4.2)]{WeightedMetricTrace}; and (iii) the
Fusco--Julin/oriented-radial-trace input to the homothetic velocity defect
\cite{WeightedMetricTrace}. None of these can be improved beyond
polynomial without altering the upstream construction.
\end{remark}

\section{Smallness regime parametrisation and the global bound}
\label{sec:smallness}

Theorem~\ref{thm:small-delta} closes the small-$\delta_T$ regime
relative to the upstream metric-trace input.
For the complementary regime $\delta_T(\Om)\ge\delta_0$, the trivial
bound $\Fr(\Om)<2$ suffices.

\begin{lemma}[Large-$\delta_T$ trivial bound]\label{lem:large-delta}
For every open $\Om\subset B_R$ with $|\Om|=\omega_n$ and
$\delta_T(\Om)\ge \delta_0(n,R,\rstar)$,
\begin{equation}\label{eq:large-delta}
\Fr(\Om)^2\le \frac{4}{\delta_0(n,R,\rstar)}\,\delta_T(\Om).
\end{equation}
\end{lemma}

\begin{proof}
$\Fr(\Om)<2$ by definition of the Fraenkel asymmetry, so
$\Fr(\Om)^2<4$. By assumption,
$1\le \delta_T(\Om)/\delta_0$, hence
$\Fr(\Om)^2\le 4\le 4\delta_T(\Om)/\delta_0$.
\end{proof}

\begin{remark}[No FMP input in the large branch]\label{rem:FMP}
The large-$\delta_T$ branch above uses only the normalization
$\Fr(\Om)<2$. No FMP quantitative isoperimetric or suboptimal
Saint--Venant estimate is needed in this block. This avoids importing
an endpoint citation whose hypotheses and constants play no role in
the route-$\delta$ transfer.
\end{remark}

\begin{proof}[Proof of Theorem~\ref{thm:main}]
For $\delta_T(\Om)\le\delta_0$, \eqref{eq:main} holds with
$C(n,R,\rstar):=C_3(n,R,\rstar)$ by
Theorem~\ref{thm:small-delta}. For $\delta_T(\Om)\ge \delta_0$,
Lemma~\ref{lem:large-delta} gives the same bound with
$C(n,R,\rstar):=4/\delta_0(n,R,\rstar)$. Taking
\begin{equation}\label{eq:C-final}
C(n,R,\rstar):=\max\!\left\{C_3(n,R,\rstar),\frac{4}{\delta_0(n,R,\rstar)}\right\}
\end{equation}
proves \eqref{eq:main} on the whole admissible class.

Sharpness of the exponent~$2$ is the standard ellipsoidal saturation:
e.g., \cite[Introduction]{BraDeP2017} shows
$\delta_T(\Om_\eps)\sim\eps^2$ for the ellipsoidal one-parameter family
$\Om_\eps$. Along this family of
volume-preserving ellipsoidal perturbations $\Om_\eps$ of $B_1$,
$\Fr(\Om_\eps)\sim\eps$ and $\delta_T(\Om_\eps)\sim\eps^2$ to second
order, so any bound $\Fr(\Om)^p\le C\delta_T(\Om)$ requires $p\ge 2$.
\end{proof}

\begin{remark}[How this slots into \textsc{GlobalAssembly}]
\label{rem:downstream}
Theorem~\ref{thm:main} is precisely the input to the bounded reduction
step of \cite[Theorem~5.1]{BDV-BoundedReduction-Plan1}, specialised to
$R=R_*(n)$. Together with the Kohler--Jobin transfer of
\cite[Theorem~6.1]{KohlerJobinTransfer-Plan1}, it yields the
sharp quantitative Faber--Krahn inequality
\[
\bigl(|\Om|/|B_1|\bigr)^{2/n}\bigl(\lambda_1(\Om)-\lambda_1(B^*)\bigr)
\ge c_{\rm FK}(n)\,\Fr(\Om)^{2},
\]
as assembled in \cite[\S6]{GlobalAssembly}.
\end{remark}

\begin{remark}[Where this block sits in the Plan~2 chain]
The boundary-layer transfer is the last \emph{geometric} step in the
route-$\delta$ chain. Upstream, \cite{LevelSetIdentity} provides the
exact deficit identity, \cite{MetricFramework} the metric first
variation, and \cite{WeightedMetricTrace} the metric trace at
$\rhad$. The present block does \emph{not} use any pointwise (R)-type
input, only the metric trace \eqref{eq:metric-trace} and the elementary
triangle inequality of \S\ref{sec:transfer}. In particular, no
$C^{1,\gamma}$ regularity of $\partial\Om$ is required; the proof is
purely finite-perimeter.
\end{remark}

\begin{thebibliography}{99}

\bibitem{LevelSetIdentity}
Plan~2 building block \textsc{LevelSetIdentity}.
Exact deficit identity for finite-perimeter superlevel sets;
Lemmas 8.1--8.3 (profile-gap measure, superlevel replacement).

\bibitem{MetricFramework}
Plan~2 building block \textsc{MetricFramework}.
Metric space $\X$, metric first variation, per-$\rho$ velocity
defect.

\bibitem{CentroidBound}
Plan~2 building block \textsc{CentroidBound}.
Auxiliary centroid estimates for the alternative space--time route.

\bibitem{SpaceTimeIdentity}
Plan~2 building block \textsc{SpaceTimeIdentity}.
Space-time $\Pi$ identity and auxiliary centroid-route diagnostics.

\bibitem{WeightedMetricTrace}
Plan~2 building block \textsc{WeightedMetricTrace}.
Integrated action $\to$ pointwise endpoint at
$\rhad=\rho_\delta$; metric trace estimate
$d_{\X}(F_{\rhad},B_1)^2\le C_*\delta_T$, uniformly on the window
$J^*=[\rho_\delta-L^*,\rho_\delta]$.

\bibitem{TrimmedVelocityRepair}
Plan~2 building block \textsc{TrimmedVelocityRepair}.
Self-truncated velocity defect replacing the pointwise lower-gradient
hypothesis.

\bibitem{GlobalAssembly}
Plan~2 building block \textsc{GlobalAssembly}.
Bounded reduction + Kohler--Jobin transfer; final sharp Faber--Krahn
inequality.

\bibitem{BDV-BoundedReduction-Plan1}
\texttt{Final/BoundedReduction.tex} (Plan~1 import).
Bounded reduction theorem at $R=R_*(n)$.

\bibitem{KohlerJobinTransfer-Plan1}
\texttt{Final/KohlerJobinTransfer.tex} (Plan~1 import).
Kohler--Jobin transfer from sharp Saint--Venant to sharp
Faber--Krahn.

\bibitem{Maggi2012}
F.~Maggi,
\emph{Sets of finite perimeter and geometric variational problems:
an introduction to geometric measure theory},
Cambridge Studies in Advanced Mathematics \textbf{135},
Cambridge University Press (2012).

\bibitem{BDV}
L.~Brasco, G.~De~Philippis, B.~Velichkov,
\emph{Faber--Krahn inequalities in sharp quantitative form},
Duke Math.\ J.\ \textbf{164} (2015), 1777--1831 (arXiv:1306.0392).

\bibitem{BraDeP2017}
L.~Brasco, G.~De~Philippis,
\emph{Spectral inequalities in quantitative form},
in: \emph{Shape Optimization and Spectral Theory},
A.~Henrot (ed.), De Gruyter Open (2017), 201--281.

\bibitem{AFP2000}
L.~Ambrosio, N.~Fusco, D.~Pallara,
\emph{Functions of bounded variation and free discontinuity problems},
Oxford Mathematical Monographs, Oxford University Press (2000).

\bibitem{FigalliMaggi2011}
A.~Figalli, F.~Maggi,
\emph{On the shape of liquid drops and crystals in the small mass regime},
Arch.\ Ration.\ Mech.\ Anal.\ \textbf{201} (2011), 143--207.

\end{thebibliography}

\end{document}
